\chapter{Spatial Sites and Lazy Domain Materialisation}

\section{Scope of this pass}

This chapter documents the next incremental step of the \texttt{Domain}
refactor.  The change is deliberately modest in algorithmic scope and strict in
architectural intent:
\begin{itemize}[nosep]
  \item keep \texttt{Domain} geometry-first;
  \item avoid materialising analysis nodes when only geometry is needed;
  \item make the different ``point-like'' entities comparable through a small
        static concept vocabulary;
  \item preserve the low-level requirements of HPC kernels: contiguous storage,
        inlining, and minimal pointer chasing.
\end{itemize}

The result is not a full removal of \texttt{Node} from the upper layers.  It is
an ontologically cleaner and computationally cheaper boundary between geometry,
analysis, and constitutive sampling.

\section{The problem of overlapping point-like concepts}

Finite-element codes naturally generate several entities that look similar
because they all carry coordinates.  However, they are not the same concept.
The overlap becomes especially visible in structural and mixed formulations:
\begin{itemize}[nosep]
  \item in Gauss--Lobatto beam integration, an endpoint may be both a nodal
        station and an integration station;
  \item in continuum formulations, a constitutive material point is collocated
        with a quadrature point but is not itself the quadrature rule;
  \item in mixed formulations, an integration site may also carry field
        unknowns, but that does not make it identical to a mesh vertex.
\end{itemize}

Therefore the correct design rule is:
\begin{center}
\emph{same spatial position does not imply same semantic type.}
\end{center}

The implementation should model collocation by \emph{binding} or
\emph{decoration}, not by collapsing all roles into a single class.

\section{A static ontology of spatial sites}

\subsection{Read-only coordinate site: \texttt{PointViewT}}

This pass introduces a small compile-time concept,
\texttt{geometry::PointViewT}.  A type satisfies it if it exposes:
\begin{itemize}[nosep]
  \item a static \texttt{dim},
  \item \texttt{coord()},
  \item \texttt{coord(i)},
  \item \texttt{coord\_ref()},
  \item \texttt{data()}.
\end{itemize}

This concept is intentionally read-only.  It captures the minimum common
geometry-facing contract needed by algorithms that only inspect coordinates.

\subsection{Mutable point: \texttt{PointT}}

\texttt{PointT} now refines \texttt{PointViewT} by adding coordinate mutation:
\begin{itemize}[nosep]
  \item \texttt{set\_coord(i, value)},
  \item \texttt{set\_coord(array)}.
\end{itemize}

This split matters because not every site should be mutable.  A constitutive
wrapper such as \texttt{MaterialPoint} should forward coordinates from the
integration point it decorates, but it should not pretend to own or mutate the
geometry itself.

\subsection{Entity-by-entity interpretation}

With the new interface split, the point-like entities can be classified as
follows:
\begin{description}[nosep]
  \item[\texttt{Point<Dim>}] mutable geometric carrier; coordinates only.
  \item[\texttt{Vertex<Dim>}] point plus stable mesh identity.
  \item[\texttt{Node<Dim, Storage>}] point plus analysis identity, DoF storage,
        and PETSc-related node metadata.
  \item[\texttt{IntegrationPoint<Dim>}] mutable quadrature site; coordinates,
        weight, and identifier.
  \item[\texttt{MaterialPoint<Policy>}] read-only constitutive decorator bound
        to an \texttt{IntegrationPoint}.
  \item[\texttt{MaterialSection<Policy>}] read-only constitutive decorator for
        structural sections, also bound to an \texttt{IntegrationPoint}.
  \item[\texttt{NodeSection<Dim>}] read-only geometric and DoF-forwarding view
        over a \texttt{Node}, enriched with section frame and section geometry.
\end{description}

The important architectural detail is that only some of these types own
coordinates.  The others forward them.

\section{Why forwarding is the right model}

\subsection{Material points and sections}

A \texttt{MaterialPoint} does not own geometry.  It owns constitutive state and
binds to the numerical sampling site where that state is evaluated.  Therefore
it now forwards:
\begin{itemize}[nosep]
  \item \texttt{coord()},
  \item \texttt{coord\_ref()},
  \item \texttt{data()},
  \item \texttt{weight()}.
\end{itemize}

The same is true for \texttt{MaterialSection}.  This makes both types valid
\texttt{PointViewT} sites without duplicating coordinate storage.

\subsection{Node sections}

\texttt{NodeSection} already bound a \texttt{Node} plus section-specific
metadata.  This pass makes the forwarding explicit for:
\begin{itemize}[nosep]
  \item coordinates (\texttt{coord}, \texttt{coord\_ref}, \texttt{data}),
  \item contiguous DoF storage (\texttt{dof\_data()}).
\end{itemize}

This is particularly useful for structural formulations, because a section
station is often treated as a special nodal view rather than as a new point
owner.

\section{Lazy node materialisation in \texttt{Domain}}

\subsection{Geometric operations should not force analysis state}

The earlier \texttt{Domain} transition already made \texttt{Vertex<Dim>} the
canonical geometric storage.  However, there were still two remaining issues:
\begin{enumerate}[nosep]
  \item \texttt{create\_boundary\_from\_plane()} still forced node-cache
        materialisation even when only geometry was needed;
  \item the non-const \texttt{node\_p(id)} path had an incorrect guard and
        could skip node-index construction.
\end{enumerate}

\subsection{Implemented fix}

This pass changes the behaviour as follows:
\begin{itemize}[nosep]
  \item \texttt{create\_boundary\_from\_plane()} is now geometry-first; it binds
        face elements to vertices immediately and binds nodes only if the node
        cache already exists;
  \item \texttt{node\_p(id)} now reliably builds the node index on demand;
  \item \texttt{Domain} exposes three explicit traceability predicates:
        \texttt{has\_materialized\_nodes()},
        \texttt{has\_vertex\_index()},
        \texttt{has\_node\_index()}.
\end{itemize}

The first predicate is important both for debugging and for architectural
tests: it allows the test suite to assert that a geometry-only workflow truly
stays geometry-only.

\subsection{Why the node cache still exists}

From a purely ontological point of view, one might want to remove
\texttt{Node} duplication entirely and replace it with a thinner attachment
structure.  That direction is still plausible, but it should not be taken
lightly.

The current element kernels, structural sections, PETSc section setup, and
element-level DoF caches still benefit from a stable \texttt{Node\&} view with:
\begin{itemize}[nosep]
  \item inline coordinate access,
  \item contiguous DoF access,
  \item zero-cost SBO-backed DoF storage policies,
  \item no proxy objects in the inner assembly path.
\end{itemize}

Keeping the node cache while making it lazy is therefore the pragmatic design
point for this stage of the refactor.

\section{Data-layout implications for HPC}

\subsection{Owned storage vs forwarded storage}

The implementation now follows a simple rule:
\begin{itemize}[nosep]
  \item \texttt{Point}, \texttt{Vertex}, \texttt{Node}, and
        \texttt{IntegrationPoint} own contiguous coordinate storage;
  \item \texttt{MaterialPoint}, \texttt{MaterialSection}, and
        \texttt{NodeSection} do not own coordinates and only forward views.
\end{itemize}

This keeps the hot data local and prevents hidden copies when decorating an
existing site with extra semantics.

\subsection{Why no generic runtime interface was added}

An obvious alternative would be an inheritance hierarchy such as:
\[
  \texttt{ISpatialSite} \leftarrow
  \{\texttt{Vertex}, \texttt{Node}, \texttt{IntegrationPoint}, \dots\}
\]
with virtual coordinate access.

That was intentionally rejected.  The reasons are both architectural and
computational:
\begin{itemize}[nosep]
  \item the library already has enough runtime polymorphism at the
        \texttt{ElementGeometry} and \texttt{Material} wrappers;
  \item coordinate access is too frequent to justify another virtual boundary;
  \item the overlap between site types is semantic, not behavioural in the
        object-oriented sense.
\end{itemize}

Compile-time concepts are the better fit here: they make the common contract
explicit without imposing a vtable or pointer-indirection tax.

\section{What this means for overlapping formulations}

The new model clarifies how to think about several important special cases:
\begin{description}[nosep]
  \item[Gauss--Lobatto beam integration] an endpoint may be represented by a
        \texttt{NodeSection} and by a \texttt{MaterialSection} bound to an
        \texttt{IntegrationPoint} at the same coordinate.  The location is
        shared; the role is not.
  \item[Continuum material points] a \texttt{MaterialPoint} is bound to an
        \texttt{IntegrationPoint}.  It is not itself the quadrature rule and
        not a mesh node.
  \item[Mixed or assumed-strain formulations] if a sampling site later carries
        unknowns, that should be modelled as an additional attachment or
        dedicated type, not by erasing the distinction between integration and
        nodal roles.
\end{description}

This is the key conceptual payoff of the refactor: the code can represent
collocation without ontological collapse.

\section{Validation}

This pass was validated by extending the test suite in two directions:
\begin{enumerate}[nosep]
  \item \texttt{domain\_vertex\_refactor} now verifies that geometry-only
        element binding and boundary creation do not materialise the node cache;
  \item \texttt{integration\_sites} verifies that
        \texttt{IntegrationPoint}, \texttt{MaterialPoint},
        \texttt{MaterialSection}, and \texttt{NodeSection} satisfy the new
        static geometry-view contract.
\end{enumerate}

This complements the earlier low-level tests on contiguous access and keeps the
refactor traceable from architecture to implementation.

\section{Interpretation}

The central lesson of this step is that the library does not need fewer
concepts; it needs cleaner boundaries between them.  The right optimisation was
not to merge \texttt{Point}, \texttt{Node}, \texttt{IntegrationPoint}, and
\texttt{Section}-like entities.  The right optimisation was:
\begin{itemize}[nosep]
  \item to expose their common geometric contract statically,
  \item to preserve ownership only where ownership is real,
  \item to keep \texttt{Domain} geometry-first for as long as possible,
  \item to materialise analysis state only when the solver path requires it.
\end{itemize}

That is a better match for both SOLID-style concept clarity and HPC-oriented
implementation discipline.
